\subsection{平展代数的微分刻画}

定理3.——设 A 是 K 上有限次数的交换代数。要使 A 是平展的，必要且充分的是 A 的 K-微分模 \(\Omega_{\mathrm{K}}(\mathrm{A})\) 约化为 0。

\begin{enumerate}
\def\labelenumi{\Alph{enumi})}
\item
  设 L 为 K 的一个代数闭包（V，第 23 页，定理 2）。要使 A 是平展的，必要且充分的是 L 上的代数 \(\mathbf{A}_{(L)}\) 可对角化（V，第30页，命题2）。进一步，A-模 \(\Omega_{\mathrm{L}}(\mathbf{A}_{(L)})\) 同构于 \(\Omega_{\mathrm{K}}(\mathbf{A}) \otimes_{\mathbf{A}} \mathbf{A}_{(L)}\) （III，第572页，命题20），从而由张量积的结合性，同构于 \(\Omega_{\mathrm{K}}(\mathbf{A}) \otimes_{\mathbf{K}} \mathbf{L}\) ；因此 \(\Omega_{\mathrm{K}}(\mathbf{A}) = 0\) 等价于 \(\Omega_{\mathrm{L}}(\mathbf{A}_{(L)}) = 0\) 。于是为证明定理3，只需考虑 K 是代数闭域的情形，并证明代数 A 可对角化当且仅当 \(\Omega_{\mathrm{K}}(\mathbf{A}) = 0\) 。
\item
  假设A可对角化；那么（V，第29页，命题1），向量空间 A 由 A 的幂等元生成。断言 \(\Omega_{\mathrm{K}}(\mathbf{A}) = 0\) 因而是以下引理的推论：
\end{enumerate}

引理2. --- 设 A 是 K 上的交换代数，且 e 是 A 的一个幂等元；则我们在 \(\Omega_{\mathbf{K}}(\mathbf{A})\) 中有 de = 0。

由关系式 \(e = e^2\) 我们推出 \(de = 2e \cdot de\) ；乘以 \(e\) 我们得到 \(e \cdot de = 2e \cdot de\) ，因此 \(e \cdot de = 0\) ，从而最终得到 \(de = 2e \cdot de = 0\) 。

\begin{enumerate}
\def\labelenumi{\Alph{enumi})}
\setcounter{enumi}{2}
\tightlist
\item
  我们首先证明两个引理：
\end{enumerate}

引理3. --- 设A是代数闭域K上的有限次交换代数，使得 \(\Omega_{\mathbf{K}}(\mathbf{A}) = 0\) 。则我们有 \(\mathfrak{m} = \mathfrak{m}^2\) （对 A 的每个极大理想 \(\mathfrak{m}\)）。

该代数 \(\mathbf{A} / \mathfrak{m}\) 是代数闭域 \(\mathbf{K}\) 的有限次扩张，从而 \([\mathbf{A} / \mathfrak{m}:\mathbf{K}] = 1\) 。因此对每个 \(a\in \mathbf{A}\) 存在唯一的标量 \(\lambda\) 使得 \(a - \lambda ,1\in \mathfrak{m}\) ；记 \(\mathrm{D}(a)\) 为 \(a - \lambda .1\) 模 \(\mathfrak{m}^2\) 的剩余类。显然 \(\mathrm{D}\) 是 \(\mathbf{A}\) 到 A-模 \(\mathfrak{m} / \mathfrak{m}^2\) 内的一个 K-导子。\(\Omega_{\mathbf{K}}(\mathbf{A})\) 的泛性质（III，第569页）及假设 \(\Omega_{\mathbf{K}}(\mathbf{A}) = 0\) 现在蕴含 \(\mathrm{D} = 0\) ，从而 \(\mathfrak{m} / \mathfrak{m}^2 = 0\) ，因此 \(\mathfrak{m} = \mathfrak{m}^2\) 。

引理4. --- 设 A 是一个交换环，且设 \(\mathfrak{a}\) 是 A 的一个有限生成理想，使得 \(\mathfrak{a} = \mathfrak{a}^2\) 。则存在一个幂等元 \(e\) （在 A 中）使得 \(\mathfrak{a} = \mathbf{A}e\) 。

设 \((a_{1},\ldots ,a_{r})\) 是理想 \(\mathfrak{a}\) 的一个生成系；由于 \(\mathfrak{a} = \mathfrak{a}^2\) ，存在元素 \(x_{ij}\) （在 \(\mathfrak{a}\) 中）使得 \(a_{i} = \sum_{j = 1}^{r}x_{ij}a_{j}\) （对 \(1\leqslant i\leqslant r\)）。我们记 M 为正方形

阶为 r 的矩阵，其元素为 \(\delta_{ij} - x_{ij}\) 并设 D 为其行列式。存在（III，第 532 页，公式 (26)）一个阶为 r 的方阵 N，其元素属于 A，使得 N · M = D · I \(_{r}\) ，由此立即得到 Da \(_{j}\) = 0 对于 \(1 \leqslant j \leqslant r\) ，因此最终 Da = 0。现在矩阵 M 模 a 与 I \(_{r}\) 相合，因此 D ≡ 1 模 a。令 e = 1 - D；则 \(e \in a\) ，且对所有 \(x \in a\) 有 ex = x。由此可知 e 是幂等元，且 a 等于 Ae。

在这些引理成立的基础上，我们通过对 A 的次数进行归纳来证明：若 K 是代数闭域且 \(\Omega_{\mathrm{K}}(\mathrm{A})=0\) ，则 A 可对角化。设 m 是 A 的一个极大理想（I，第 104 页）。由引理 3 和引理 4，存在一个幂等元 e 使得 m = Ae；我们已经看到 A/m 在 K 上的次数为 1。因此 A 是理想 \(\mathfrak{a}=(1-e)\mathbf{A}\) 与 m 的直和，并且我们有 \([\mathfrak{a}:K]=1\) ，因此 A 同构于 \(K\times A/a\) 。由于 \(\Omega_{\mathrm{K}}(\mathrm{A}/\mathfrak{a})\) 同构于 \(\Omega_{\mathrm{K}}(\mathrm{A})\) 的一个商（III，第 573 页，命题 22），它为零，且归纳假设表明 A/ \(\mathfrak{a}\) 是可对角化的。这进而表明 A 是可对角化的。

